\subsection{LeNet-5: The First Famous CNN}\label{sec:lenet-5-the-first-famous-cnn}

\subsubsection{Historical Motivation}\label{sec:historical-motivation}

Let's move on the first concerete CNN architecture: LeNet-5, which was proposed by \emph{Yann LeCun} in 1998 \cite{lecun1998gradient}. It is one of the first successful convolutional neural networks. It was designed for \textbf{handwritten-digit recognition}: given a small grayscale image, the network predicts one of the ten digit classes.

\highspace
\begin{flushleft}
    \textcolor{Red2}{\faIcon{exclamation-triangle} \textbf{The problem: handwritten-digit recognition}}
\end{flushleft}
\definition{LeNet-5} was introduced by \emph{Yann LeCun} and collaborators for \textbf{handwritten digit recognition}. The input consists of very small \textbf{grayscale images}, and the goal is straightforward:
\begin{equation*}
    \text{image of a digit} \to \text{digit class (0, 1, 2, 3, 4, 5, 6, 7, 8, or 9)}
\end{equation*}
For example, an image containing a handwritten digit `3' should be classified as the digit class `3'. The recognition of single handwritten digits can then be used to recognize complete handwritten numbers as a \textbf{sequence of digits} (e.g., a handwritten number `123' can be recognized as the sequence of digits `1', `2', and `3').

\highspace
So, conceptually, this is exactly the image-classification problem we have been discussing in the previous sections:
\begin{equation*}
    I \in \mathbb{R}^{H \times W \times 1} \to \text{CNN} \to \text{10 classes (0, 1, 2, 3, 4, 5, 6, 7, 8, or 9)}
\end{equation*}
The $1$ in the input shape indicates that the input images are \textbf{gray-scale} (single channel).

\highspace
\textcolor{Green3}{\faIcon{question-circle} \textbf{Why not simply use an Multi-Layer Perceptron (MLP)?}} We already know that this is possible, but it is \textbf{inefficient}. We could flatten the image immediately: $H \times W \to HW$, and feed every pixel directly into an MLP. Technically, \textbf{this works}, there is nothing mathematically preventing us from doing it. But we should \textbf{not} treat every pixel as an independent input feature because images are \textbf{highly spatially correlated}.

\highspace
\textcolor{Green3}{\faIcon{check-circle} \textbf{Why is this more efficient?}} Imagine connecting the $32 \times 32 = 1024$ pixels directly to just 84 dense neurons: $1024 \times 84 = 86'016$ weights already. And every neuron learns completely separate weights for every pixel. For now, it is sufficient to know that the first convolutional layer of the LeNet-5 architecture contains six $5 \times 5$ filters. Since the image is gray-scale $C_{\text{in}} = 1$ and thus:
\begin{equation*}
    \text{\# params} = 6 \left(5 \cdot 5 \cdot 1 + 1\right) = 156
\end{equation*}
Those 156 parameters are enough to apply six learned feature detectors \textbf{everywhere in the image}. Thus:
\begin{itemize}
    \item MLP: different weights for different pixel connections.
    \item CNN: small local filters reused across spatial positions.
\end{itemize}